% Drawing 74/02-1 -- folder 74, batch 02, archive page 21 (PDF page 22).
% Transcription: un prisme hexagonal (ou un cube tronqué) en perspective ; au-dessous, un graphe en éventail : un sommet relié par des arcs à cinq sommets alignés, et un petit trièdre de trois segments issus d'un point.
% Kind: polyhedron sketch, fan graph, trihedron. Method 2 (schematic redraw from the facsimile). Status: ok.
% Notes: Drawn over a printed examination form (the sheet is upside down; its dotted rules and typed headings are not drawn). Top: a hexagonal outline with an inner vertex joined to three alternate corners (a cube in perspective, overinked); middle: an apex joined to five aligned vertices, the first two of which are also joined to a lower point; bottom: a trihedron, three segments from one point.
% Coordinates: 1 unit = 1 mm = 1/100 of the crop width (bbox in page fractions: [0.12, 0.43, 0.42, 0.68]).
% Redrawn by Claude (claude-opus-5-5), issue #73, 2026-11-07; not checked by a human; not compiled (no TeX in the session).
% Rights: reproduces part of the work -- behind RENDER_TRANSCRIPTS, never published by itself.
\documentclass[tikz,border=3pt]{standalone}
\usepackage{amsmath,amssymb}
\usetikzlibrary{patterns,arrows.meta}
\tikzset{pencil/.style={black!45}, struck/.style={black!35}, mark/.style={circle,draw,inner sep=1.3pt}, pt/.style={circle,fill,inner sep=0.9pt}, opt/.style={circle,draw,fill=white,inner sep=1.1pt}}
\begin{document}
\begin{tikzpicture}[x=1mm,y=1mm,line cap=round,line join=round,font=\scriptsize]
\draw (45,-24) -- (55,-10) -- (70,-15) -- (73,-28) -- (64,-35) -- (50,-36) -- cycle;
\draw (45,-24) -- (59,-25) -- (70,-15);
\draw (59,-25) -- (64,-35);
\draw (40,-43) -- (22,-66);
\draw (40,-43) -- (40,-66);
\draw (40,-43) -- (58,-65);
\draw (40,-43) -- (70,-66);
\draw (40,-43) -- (82,-63);
\node[pt] at (22,-66) {};
\node[pt] at (40,-66) {};
\node[pt] at (58,-65) {};
\node[pt] at (70,-66) {};
\node[pt] at (82,-63) {};
\draw (22,-66) -- (38,-84) -- (40,-66);
\draw (80,-108) -- (67,-95);
\draw (80,-108) -- (80,-93);
\draw (80,-108) -- (93,-93);
\end{tikzpicture}
\end{document}
